\begin{abstract}
O dimensionamento de pontes de madeira envolve múltiplas variáveis de projeto e critérios normativos conflitantes, sendo tradicionalmente conduzido por procedimentos iterativos dependentes da experiência do projetista.
Este trabalho apresenta uma abordagem de projeto inteligente para pontes de madeira roliça de estradas vicinais, na qual a modelagem paramétrica dos elementos estruturais é integralmente acoplada a uma rotina de otimização robusta multiobjetivo.
O problema é formulado com duas funções objetivo antagônicas, ambas de minimização: o volume total de madeira consumido e a utilização do limite de serviço de deslocamento, $\delta_{tot}/(L/350)$, sujeitas às restrições de estado limite último e de serviço prescritas pela ABNT NBR 7190:2022.
A resolução emprega o Algoritmo Genético de Ordenação Não Dominada II (NSGA-II), totalmente acoplado ao modelo determinístico de verificação da longarina e do tabuleiro.
A robustez é incorporada de forma nativa ao processo de busca: cada solução candidata é avaliada considerando um desvio de $\rho = 5\%$ aplicado aos valores nominais das variáveis de projeto, e as funções objetivo e restrições são agregadas por múltiplas checagens, de modo que apenas soluções tolerantes a variações permaneçam na fronteira eficiente.
A metodologia foi implementada em uma plataforma computacional de acesso livre e aplicada a uma matriz de 20 configurações de ponte vicinal de pista simples, combinando quatro vãos teóricos ($L = 3{,}0$~m, $4{,}0$~m, $5{,}0$~m e $6{,}0$~m) e as cinco classes de resistência de madeiras de florestas nativas da ABNT NBR 7190-3 (D20 a D60), todas submetidas ao trem tipo TB-240.
As vinte células convergiram para soluções viáveis, com consumo de madeira entre 1{,}75 e 9{,}98~m\textsuperscript{3}, equivalentes a 0{,}129 e 0{,}370~m\textsuperscript{3} por metro quadrado de tabuleiro.
O consumo cresce com o vão segundo uma lei de potência cujo expoente varia pouco entre classes, entre 1{,}64 e 1{,}73, e o primeiro degrau de classe é o mais rentável: passar de D20 para D30 reduz o consumo em cerca de 17 a 21\% em qualquer um dos quatro vãos, cerca de 41\% a 48\% de toda a economia disponível até D60, ao passo que os degraus seguintes rendem, individualmente, entre 4 e 12\%.
Em todas as vinte células a solução econômica opera a no máximo 1{,}4\% de folga do limite normativo da verificação governante, sem saturar o piso dimensional do domínio de busca em nenhuma delas.
Contrariando parcialmente a expectativa corrente, a flexão governa a totalidade das vinte células, quinze pelo tabuleiro e cinco pela longarina, sem nenhuma célula governada pela flecha, ainda que a flecha da longarina chegue extremamente perto do limite em três células de vão e classe mais severos, um indício de que a migração de resistência para rigidez começa a se anunciar dentro da própria faixa de vãos estudada.
A célula de referência foi ainda reotimizada para uma grade de cinco níveis de desvio, $\rho \in \{0\%,\ 5\%,\ 10\%,\ 15\%,\ 20\%\}$: ao contrário do que uma versão anterior deste estudo reportou, esta rodada detectou um custo de robustez real e monotônico, de $+6{,}0\%$ em $\rho = 5\%$ a $+26{,}9\%$ em $\rho = 20\%$ em relação ao caso determinístico, com taxa marginal entre 1{,}2\% e 1{,}4\% de volume adicional por ponto percentual de $\rho$.
Os resultados da varredura foram ainda condensados em ábacos gráficos de anteprojeto e em uma equação fechada de pré-dimensionamento do volume, ajustada por regressão sobre as vinte células com $R^2 = 0{,}999$ e erro médio de 1{,}3\%.
A plataforma, os ábacos e a equação de anteprojeto dela derivados contribuem para a automação, a racionalização e a modernização do projeto de pontes de madeira em vias vicinais.
\end{abstract}


\textbf{Palavras-chave}: pontes de madeira, estradas vicinais, projeto inteligente, otimização multiobjetivo, NSGA-II, otimização robusta, NBR 7190.
